\section{Evaluation}
\label{sec:evaluation}
% ============================================================
% MARKING CRITERION: "Evaluation" [20 marks]
%   - Evaluation methods and metrics discussed thoroughly
%   - Appropriate evaluation plan
%   - Evaluation method clear and sensible
%   - Results / conclusions presented clearly

\guidance{%
  MARKING: This section is worth 20 marks. The marker will check:
  (1) that evaluation methods and metrics are discussed thoroughly,
  (2) that you had an appropriate evaluation plan, (3) that the method
  is clear and sensible, and (4) that results and conclusions are
  presented clearly. Use tables and figures to present results.
  Interpret your results --- do not just state numbers.}

Briefly introduce your evaluation strategy, the scenario(s) tested, and any datasets used.

\subsection{Intrinsic Evaluation}
\label{sec:intrinsic}

\guidance{%
  Describe the intrinsic evaluation: what was measured, what data was
  used (e.g.\ dataset, number of test items), and what metrics were
  applied. Define each metric with a formula if appropriate.}

% Define your metrics:
% \begin{equation}
%   \text{MetricName} = \frac{\text{Numerator}}{\text{Denominator}}
%   \label{eq:metric}
% \end{equation}

\subsubsection{Results}

\guidance{%
  Present intrinsic results clearly in a table. Interpret them ---
  what do they tell you about your system's performance? Compare with
  baselines or prior work where possible.}

\begin{table}[H]
\centering
\small
\begin{tabular}{lcc}
\toprule
\textbf{Metric} & \textbf{Condition A} & \textbf{Condition B} \\
\midrule
Metric 1  & 0.000 & 0.000 \\
Metric 2  & 0.000 & 0.000 \\
Metric 3  & 0.000 & 0.000 \\
\bottomrule
\end{tabular}
\caption{Intrinsic evaluation results.}
\label{tab:intrinsic_results}
\end{table}

% Interpret the results here.


\subsection{Extrinsic Evaluation}
\label{sec:extrinsic}

\guidance{%
  Describe your extrinsic evaluation plan: participant recruitment,
  study procedure, ethical approval, and any measures taken to ensure
  data quality (e.g.\ attention checks, exclusion criteria).}

\subsubsection{Task Performance}

\guidance{%
  Report objective task-level results: task success rate, number of
  turns, slot-filling accuracy, dialogue efficiency, etc.}

% Write task performance results here.

\subsubsection{Behavioural Analysis}

\guidance{%
  If applicable, analyse how the system handled specific user behaviours
  (e.g.\ pauses, self-corrections, interruptions). Present in a table.}

\begin{table}[H]
\centering
\small
\begin{tabular}{lcc}
\toprule
\textbf{Behaviour} & \textbf{Occurrences} & \textbf{Success Rate} \\
\midrule
Behaviour A & 0 & 0\% \\
Behaviour B & 0 & 0\% \\
Behaviour C & 0 & 0\% \\
\bottomrule
\end{tabular}
\caption{Success rates for handling different behaviours.}
\label{tab:behaviour_results}
\end{table}

\subsubsection{Subjective Evaluation}

\guidance{%
  Report questionnaire results. Describe the instrument (e.g.\
  Godspeed, Likert scales), validity checks (KMO, Cronbach's alpha),
  and statistical tests used. Present frequencies and averages with
  figures if possible.}

% Write subjective evaluation results here.

\subsection{Discussion of Results}

\guidance{%
  Tie all evaluation results together. What do the intrinsic and
  extrinsic results tell you collectively? Were there any surprising
  findings? What are the limitations of your evaluation? This is
  where you demonstrate clear and sensible interpretation --- a key
  part of the marking criterion.}

% Write discussion here.
